\chapter{Predictor Hyperparameter Searches}\label{app:predictor_search}

Tables~\ref{tab:search_spectral_mlp}--\ref{tab:search_waveform_cnn} list the \ac{mlp} and \ac{cnn} trials from Section~\ref{sec:offline_predictor_benchmark} in trial order. The first 15 trials use randomly chosen hyperparameter sets while the rest are selected using the tree-structured Parzen estimator by Optune.
Bold rows mark the configuration with the lowest score in each search.
Figure~\ref{fig:predictor_search_loss_curves} shows the loss curves for the selected fits, and Figure~\ref{fig:predictor_search_dmd_heatmaps} summarizes the spectral \ac{dmdc} grid.

% Keep the loss curves with the appendix introduction.
\begin{figure}[!htbp] % latex-lint:ignore=FLOAT-01
  \centering
  \includegraphics{figures/08_appendix/predictor_search_loss_curves.pdf}
  \caption{Training and validation losses for the selected \ac{mlp} and \ac{cnn} predictors.
  The jumps in training loss occur when the curriculum lengthens the forecast scored by the observation-space \ac{mse}, adding errors from later steps.
  This horizon grows from one step to 17 frames for spectral log-power or 12 samples for the waveform.
  The spectral term in the waveform joint loss always spans 17 frames, and validation uses the final horizons throughout.
  The dashed line marks the first epoch eligible for checkpoint selection, when the full training horizon is scored and all terms are active.
  The circle marks the eligible epoch with the lowest logged validation loss.
  The spectral \ac{cnn} sweep ranked the resulting checkpoints by a separate recursive log-power \ac{mse} evaluation.}\label{fig:predictor_search_loss_curves}
\end{figure}

% Keep the heatmaps above the first trial table.
\begin{figure}[!t] % latex-lint:ignore=FLOAT-01
  \centering
  \includegraphics{figures/08_appendix/predictor_search_dmd_heatmaps.pdf}
  \caption{Validation loss across 140 spectral \ac{dmdc} fits.
  For each \ac{eeg} history \(\nhist{o}\) and retained fraction of the squared singular-value sum, the left panel shows the lowest validation log-power \ac{mse} across five ridge values \(\lambda\).
  The circle marks the minimum of \(1.6197\) at four frames, a retained fraction of \(0.95\), and \(\lambda=0.01\).
  The right panel shows the range across ridge values in units of \(10^{-7}\), indicating little effect from ridge at fixed history and retained fraction.}\label{fig:predictor_search_dmd_heatmaps}
\end{figure}

\begin{table}[!t] % latex-lint:ignore=FLOAT-01
  \centering
  \small
  \caption{Spectral \ac{mlp} trials scored by validation log-power \ac{mse} over 17 forecast frames.
  History counts input frames.
  Trials 0--24 predict the increment from the latest frame.
  Trial 25 uses the architecture of trial 23 to predict the next frame directly and gives the lowest loss.}\label{tab:search_spectral_mlp}
  % Generated by predictor_search_appendix.py.
\begin{tabular}{@{}rrrcllr@{}}
\toprule
Trial & History & Width & Layers & Activation & Output & Score \\
\midrule
0 & 3 & 128 & 3 & softplus & Residual & 2.2737 \\
1 & 8 & 512 & 3 & tanh & Residual & 1.9927 \\
2 & 3 & 128 & 3 & softplus & Residual & 2.2775 \\
3 & 4 & 256 & 3 & softplus & Residual & 2.0919 \\
4 & 2 & 768 & 3 & tanh & Residual & 1.7706 \\
5 & 4 & 256 & 3 & tanh & Residual & 2.0169 \\
6 & 7 & 256 & 3 & softplus & Residual & 2.1327 \\
7 & 7 & 256 & 2 & tanh & Residual & 2.0189 \\
8 & 2 & 512 & 2 & tanh & Residual & 1.7830 \\
9 & 6 & 128 & 1 & tanh & Residual & 2.1196 \\
10 & 9 & 768 & 3 & tanh & Residual & 2.0016 \\
11 & 3 & 128 & 3 & softplus & Residual & 2.2805 \\
12 & 4 & 256 & 1 & softplus & Residual & 1.9376 \\
13 & 4 & 256 & 3 & softplus & Residual & 2.0834 \\
14 & 5 & 256 & 3 & tanh & Residual & 2.0218 \\
15 & 1 & 512 & 2 & tanh & Residual & 1.7546 \\
16 & 3 & 768 & 2 & tanh & Residual & 1.7688 \\
17 & 4 & 768 & 2 & tanh & Residual & 1.7974 \\
18 & 2 & 768 & 2 & softplus & Residual & 1.7476 \\
19 & 2 & 768 & 1 & softplus & Residual & 1.6956 \\
20 & 2 & 768 & 1 & softplus & Residual & 1.6763 \\
21 & 1 & 768 & 1 & softplus & Residual & 1.6807 \\
22 & 1 & 768 & 1 & softplus & Residual & 1.6851 \\
23 & 1 & 1024 & 1 & softplus & Residual & 1.6541 \\
24 & 1 & 1024 & 1 & softplus & Residual & 1.6672 \\
\textbf{25} & \textbf{1} & \textbf{1024} & \textbf{1} & \textbf{softplus} & \textbf{Direct} & \textbf{1.6340} \\
\bottomrule
\end{tabular}

\end{table}

\begin{table}
  \centering
  \small
  \caption{Spectral \ac{cnn} trials scored by recursive validation log-power \ac{mse} over 17 frames.
  \ac{eeg} history is in frames and kernels are given as frames by frequency bins.
  The Channels column gives the width of the first convolutional layer, and Taper is the factor applied before rounding down the width of each subsequent layer.
  The Head column gives the hidden width or its sizing rule, with Affine indicating no hidden layer.
  Trial 25 is the baseline, and trial 26 uses the selected architecture of trial 22 with a kernel spanning one frame instead of three.}\label{tab:search_spectral_cnn}
  % Generated by predictor_search_appendix.py.
\begin{tabular}{@{}rrrcrcllr@{}}
\toprule
Trial & History & Channels & Layers & Taper & Kernel & Head & Activation & Score \\
\midrule
0 & 8 & 28 & 1 & --- & \(3\times5\) & Affine & tanh & 1.7127 \\
1 & 14 & 26 & 2 & 0.6618 & \(3\times3\) & 133 & softplus & 1.6348 \\
2 & 11 & 15 & 1 & --- & \(3\times1\) & 160 & tanh & 1.7043 \\
3 & 17 & 21 & 2 & 0.5469 & \(3\times3\) & 134 & tanh & 1.6775 \\
4 & 15 & 12 & 1 & --- & \(3\times3\) & Affine & tanh & 1.6709 \\
5 & 17 & 21 & 2 & 0.7664 & \(3\times3\) & Affine & tanh & 1.7187 \\
6 & 7 & 21 & 1 & --- & \(3\times5\) & Affine & softplus & 1.6616 \\
7 & 8 & 24 & 1 & --- & \(3\times5\) & 34 & softplus & 1.6766 \\
8 & 14 & 14 & 2 & 0.5975 & \(3\times3\) & Affine & tanh & 1.6913 \\
9 & 6 & 18 & 2 & 0.8087 & \(3\times5\) & Affine & tanh & 1.6538 \\
10 & 12 & 25 & 2 & 0.7829 & \(3\times1\) & Affine & softplus & 1.6362 \\
11 & 7 & 30 & 2 & 0.9842 & \(3\times5\) & 259 & softplus & 1.6364 \\
12 & 10 & 13 & 2 & 0.9443 & \(3\times5\) & 123 & tanh & 1.6619 \\
13 & 14 & 18 & 1 & --- & \(3\times3\) & 156 & tanh & 1.7170 \\
14 & 9 & 18 & 1 & --- & \(3\times1\) & Affine & softplus & 1.6677 \\
15 & 11 & 24 & 2 & 0.7136 & \(3\times3\) & Affine & softplus & 1.6443 \\
16 & 12 & 25 & 2 & 0.6722 & \(3\times3\) & 627 & softplus & 1.7265 \\
17 & 14 & 27 & 2 & 0.8519 & \(3\times1\) & Affine & softplus & 1.6722 \\
18 & 15 & 32 & 2 & 0.6475 & \(3\times3\) & 34 & softplus & Pruned \\
19 & 15 & 14 & 2 & 0.5071 & \(3\times1\) & 69 & softplus & 1.6188 \\
20 & 14 & 32 & 1 & --- & \(3\times3\) & 61 & softplus & 1.6902 \\
21 & 14 & 17 & 2 & 0.5047 & \(3\times1\) & 72 & softplus & 1.6195 \\
\textbf{22} & \textbf{16} & \textbf{16} & \textbf{2} & \textbf{0.5015} & \textbf{\(3\times1\)} & \textbf{78} & \textbf{softplus} & \textbf{1.5996} \\
23 & 14 & 16 & 2 & 0.5043 & \(3\times1\) & 67 & softplus & 1.6259 \\
24 & 17 & 10 & 2 & 0.5077 & \(3\times1\) & 60 & softplus & 1.6248 \\
25 & 9 & 32 & 2 & 0.5000 & \(3\times3\) & Last conv. & tanh & 1.7149 \\
26 & 16 & 16 & 2 & 0.5015 & \(1\times1\) & 78 & softplus & 1.6124 \\
\bottomrule
\end{tabular}

\end{table}

\begin{table}
  \centering
  \small
  \caption{Waveform \ac{mlp} trials scored by joint validation loss at the selected epoch.
  The score adds waveform \ac{mse} over 12 samples and 15 times log-power \ac{mse} over 17 frames.
  \ac{eeg} history is in samples and current history is in milliseconds, excluding the candidate current.
  Residual output predicts the increment from the latest waveform sample, and direct output predicts the next sample itself.}\label{tab:search_waveform_mlp}
  % Generated by predictor_search_appendix.py.
\begin{tabular}{@{}rrrrcllr@{}}
\toprule
Trial & History & Current & Width & Layers & Activation & Output & Score \\
\midrule
0 & 40 & 240 & 64 & 3 & softplus & Direct & 72.2568 \\
1 & 11 & 80 & 256 & 3 & softplus & Residual & 68.8750 \\
2 & 28 & 480 & 64 & 1 & tanh & Direct & 72.1262 \\
3 & 6 & 240 & 256 & 2 & tanh & Residual & 67.0733 \\
4 & 28 & 240 & 128 & 2 & tanh & Residual & 68.8017 \\
5 & 30 & 80 & 64 & 1 & softplus & Residual & 73.2009 \\
6 & 3 & 80 & 256 & 1 & softplus & Direct & 77.1582 \\
7 & 28 & 480 & 128 & 1 & tanh & Residual & 71.0344 \\
8 & 23 & 480 & 256 & 1 & tanh & Direct & 69.5348 \\
9 & 45 & 480 & 128 & 2 & softplus & Residual & 72.2522 \\
10 & 21 & 80 & 128 & 1 & softplus & Direct & 67.3837 \\
11 & 25 & 480 & 128 & 3 & tanh & Direct & 68.3325 \\
\textbf{12} & \textbf{13} & \textbf{240} & \textbf{256} & \textbf{2} & \textbf{tanh} & \textbf{Residual} & \textbf{63.1484} \\
13 & 23 & 80 & 256 & 1 & softplus & Residual & 68.3970 \\
14 & 13 & 480 & 64 & 1 & tanh & Direct & 71.9878 \\
15 & 8 & 240 & 256 & 2 & tanh & Residual & 65.7149 \\
16 & 19 & 240 & 256 & 2 & tanh & Residual & 65.9439 \\
17 & 16 & 240 & 64 & 3 & tanh & Residual & 70.2486 \\
18 & 32 & 80 & 256 & 2 & tanh & Residual & 69.8074 \\
19 & 8 & 480 & 256 & 2 & softplus & Residual & 69.3971 \\
20 & 9 & 240 & 128 & 1 & tanh & Residual & 71.2217 \\
21 & 18 & 240 & 256 & 2 & tanh & Residual & 65.4048 \\
22 & 9 & 240 & 256 & 3 & tanh & Direct & 64.6212 \\
23 & 4 & 240 & 256 & 3 & tanh & Direct & 68.1040 \\
24 & 17 & 240 & 256 & 3 & tanh & Direct & 65.5611 \\
\bottomrule
\end{tabular}

\end{table}

\begin{table}
  \centering
  \small
  \caption{Waveform \ac{cnn} trials scored by the same joint validation loss as in Table~\ref{tab:search_waveform_mlp}.
  \ac{eeg} history and kernel length are in samples and current history is in milliseconds, excluding the candidate current.
  Channels and taper describe the convolutional layers as in Table~\ref{tab:search_spectral_cnn}.
  A hidden prediction head takes its width from either the last convolutional layer or the geometric mean of the head's input and output dimensions.
  The selected fit uses one convolutional layer with a three-sample kernel and the geometric-mean head.}\label{tab:search_waveform_cnn}
  % Generated by predictor_search_appendix.py.
\begin{tabular}{@{}rrrrcrrllr@{}}
\toprule
Trial & History & Current & Channels & Layers & Taper & Kernel & Head & Activation & Score \\
\midrule
0 & 62 & 240 & 22 & 1 & --- & 4 & Geom. mean & tanh & 77.0492 \\
1 & 33 & 480 & 19 & 1 & --- & 3 & Geom. mean & softplus & 73.5247 \\
2 & 37 & 240 & 16 & 2 & 0.5583 & 3 & Last conv. & tanh & 85.9020 \\
3 & 41 & 80 & 28 & 1 & --- & 2 & Affine & tanh & 76.7570 \\
4 & 71 & 240 & 21 & 3 & 0.9281 & 4 & Affine & tanh & 78.2708 \\
5 & 47 & 240 & 11 & 3 & 0.9891 & 3 & Geom. mean & tanh & 78.9462 \\
6 & 50 & 480 & 9 & 2 & 0.7593 & 3 & Last conv. & tanh & 88.6063 \\
7 & 13 & 80 & 12 & 1 & --- & 4 & Geom. mean & softplus & 78.8224 \\
8 & 51 & 240 & 25 & 2 & 0.5336 & 5 & Affine & softplus & 78.4380 \\
9 & 69 & 80 & 22 & 1 & --- & 4 & Affine & softplus & 77.4275 \\
10 & 41 & 240 & 13 & 3 & 0.6657 & 2 & Geom. mean & tanh & 82.4461 \\
11 & 41 & 240 & 26 & 1 & --- & 2 & Last conv. & softplus & Failed \\
12 & 62 & 240 & 22 & 1 & --- & 4 & Geom. mean & tanh & 76.2730 \\
13 & 33 & 480 & 19 & 1 & --- & 3 & Geom. mean & softplus & 74.0109 \\
14 & 37 & 240 & 16 & 2 & 0.5583 & 3 & Last conv. & tanh & 85.0890 \\
15 & 41 & 80 & 28 & 1 & --- & 2 & Affine & tanh & 77.4713 \\
16 & 31 & 480 & 18 & 1 & --- & 3 & Geom. mean & softplus & 75.2389 \\
17 & 26 & 480 & 16 & 1 & --- & 3 & Affine & softplus & 77.3876 \\
18 & 19 & 480 & 25 & 1 & --- & 4 & Geom. mean & softplus & 72.6356 \\
19 & 43 & 480 & 23 & 1 & --- & 4 & Geom. mean & softplus & 73.5362 \\
20 & 18 & 480 & 27 & 1 & --- & 3 & Geom. mean & softplus & 72.6933 \\
21 & 26 & 480 & 29 & 1 & --- & 2 & Last conv. & softplus & 76.5267 \\
22 & 18 & 480 & 25 & 1 & --- & 3 & Geom. mean & softplus & 72.7234 \\
23 & 21 & 480 & 23 & 1 & --- & 5 & Geom. mean & softplus & 72.7319 \\
24 & 13 & 480 & 21 & 2 & 0.8265 & 3 & Geom. mean & softplus & 79.2463 \\
\textbf{25} & \textbf{16} & \textbf{240} & \textbf{24} & \textbf{1} & \textbf{---} & \textbf{3} & \textbf{Geom. mean} & \textbf{softplus} & \textbf{72.4256} \\
\bottomrule
\end{tabular}

\end{table}

\chapter{Linear Prediction of Waveform and Spectral Observables}\label{app:dmdc_comparison}

The Koopman perspective in Section~\ref{sec:linear_dmd_predictor} motivates fitting linear dynamics to nonlinear functions of the measurements.
An exact finite-dimensional description requires closure of the chosen observables under evolution, whereas the regressions here test prediction over a finite horizon~\cite{Brunton2016,Williams2015}. % what is that supposed to mean? Not a coherent argument. % There are also emprical and approximative koopman approaches. So whats the point? 
Korda and Mezi\'{c}~\cite{Korda2018} use lifted predictors for \ac{mpc} and show how observables can enter both the cost and constraints. % this citation style: "Korda and Mezi\'{c}~\cite{Korda2018}" is not wanted. Citation comes at the end of the sentence without authors names. % so? 
Here the lift operates on \ac{eeg} histories and the cost penalizes predicted log-power above a healthy envelope.
The controller can therefore use the spectral prediction directly, without reconstructing waveform phase or the underlying neural state. % So? why is that important?
% this whole paragraph is just a series of statements. But there is no real argumentation that the reader can follow.

\section{Comparison with Nonlinear Predictors}

We measure the difference in forecast error between \ac{dmdc} and its nonlinear peers within each prediction target. % its the validation error
The comparison uses the first \SI{240}{\milli\second}, where the waveform prediction retains useful phase information.
Within each target every predictor starts from the same validation origins, receives the recorded future currents, and predicts recursively without further measurements.
For waveform prediction, the neural fits use waveform \ac{mse} alone and are scored at 12 sample endpoints on \num{4120} origins from 40 regenerated validation trajectories.
The spectral errors use four frame endpoints on \num{6400} origins from the 40 validation trajectories with \SI{60}{\milli\second} current holds.

% Data: waveform_dmd/pure_mse_comparison/matched_cnn_240ms.json and waveform_matched_geometry/offline/lookahead.csv (first four rows).
\begin{table}
  \centering
  \caption{Mean recursive prediction errors through \SI{240}{\milli\second}. Each row uses the units of its prediction target. The gap is the relative increase in \ac{mse} of \ac{dmdc} over the best neural fit.}\label{tab:dmdc_representation_gap}
  \begin{tabular}{lrrrr}
    \toprule
    Prediction target & \ac{dmdc} & \ac{mlp} & \ac{cnn} & Relative gap \\
    \midrule
    Waveform             & 39.07 & 29.54 & 35.49 & \(+32.3\%\) \\
    Spectral log-power   & 1.308 & 1.368 & 1.393 & \(-4.4\%\)  \\
    \bottomrule
  \end{tabular}
\end{table}

The nonlinear \ac{mlp} improved short-horizon waveform prediction, whereas the spectral \ac{dmdc} matched or improved on both nonlinear fits over the same physical duration (Table~\ref{tab:dmdc_representation_gap}).
The result supports a useful affine approximation to the evolution of the measured spectra in these trajectories.
The increase in feature dimension accompanies a loss of phase information and temporal averaging from the overlapping spectral windows, both of which can reduce the variation that the predictor has to explain.
The history lengths, excitation schedules, and selection criteria also differ, including selection of the waveform \ac{dmdc} fit at \SI{1.02}{\second}.
The experiment therefore leaves open how much of the smaller spectral gap comes from any one of these changes, and whether the chosen observables are closed under evolution.

\section{Quadratic Programming Formulation}

With a fixed initial history, the affine \ac{dmdc} recurrence makes the stacked predicted log-power \(\mathbf{o}\) affine in the stacked physical currents \(\mathbf{u}\).
Write \(\mathbf{o}=\mathbf{b}+\mathbf{G}\mathbf{u}\), where \(\mathbf{b}\) is the zero-current rollout from that history and \(\mathbf{G}\) is constant for a fixed predictor and control horizon.
Let \(\mathbf{e}\) be the healthy envelope repeated over the prediction endpoints.
An epigraph reformulation of the positive part in~\eqref{eq:mpc_spectral_cost} introduces one auxiliary variable per predicted spectral value.
\begin{equation} % this equation is not great. I dont feel like it fits into the style of the thesis. It reused symbols. Its lazy by just adding text. there is already an exaple of the epigraph reformulation in another chapter.
    \begin{aligned}
        \min_{\mathbf{u},\,\mathbf{s}}\quad
        & \mathbf{u}^{\mathsf T}\mathbf{R}\mathbf{u}
          +\mathbf{s}^{\mathsf T}\mathbf{W}\mathbf{s} \\
        \text{subject to}\quad
        & \mathbf{s}\geq\mathbf{0},\qquad
          \mathbf{s}\geq\mathbf{b}+\mathbf{G}\mathbf{u}-\mathbf{e}, \\
        & \text{the current bounds and Kirchhoff equalities in~\eqref{eq:mpc_current_constraints}.}
    \end{aligned}
    \label{eq:dmdc_spectral_hinge_qp}
\end{equation}
The diagonal matrices \(\mathbf{R}\) and \(\mathbf{W}\) carry the nonnegative weights and normalization for \(\mathcal J_{L_2}\) and \(\mathcal J_{\mathrm{spec}}\), respectively.
For positive spectral cost weights, the minimizing slack is the elementwise positive part of \(\mathbf{o}-\mathbf{e}\), so the reformulation preserves the original objective after minimizing over the slacks.
The problem has affine constraints and a positive semidefinite quadratic objective, giving a convex quadratic program for linear \ac{mpc}.
The waveform \ac{dmdc} recurrence is affine as well, but the same spectral cost on a waveform rollout requires evaluating windowed power and the logarithm inside the optimization.
Those operations make the envelope residual nonlinear in the planned currents.

The condensed implementation solves~\eqref{eq:dmdc_spectral_hinge_qp} with OSQP for continuous, independently chosen currents, equal positive stage and terminal spectral cost weights, and zero \(L_1\) effort.
The auxiliary variables here price spectral excess above the healthy envelope, while those in~\eqref{eq:mpc_l1_epigraph} price absolute current.
Eliminating the predictor states leaves 30 current variables and \num{7440} spectral slack variables at the ten-frame control horizon used here.
The solver reuses the fixed quadratic and constraint matrices and updates the constraint bounds from the new zero-current rollout at each replan.

\section{Closed-Loop Performance and Solve Time}

On ten paired reserved simulation seeds with a \SI{0.60}{\second} control horizon, the spectral and waveform \ac{dmdc} controllers gave mean seizure burdens of \(5.00\%\) and \(5.30\%\), respectively, with mean delivered charges of \SI{25.61}{\milli\ampere\second} and \SI{20.59}{\milli\ampere\second}.
Both used IPOPT under their respective solver settings, with per-seed median solve times averaging \SI{0.967}{\second} for spectral \ac{dmdc} and \SI{0.715}{\second} for waveform \ac{dmdc}.
Both replanned every \SI{60}{\milli\second}.
The waveform controller applied three independently optimized \SI{20}{\milli\second} currents in that interval, while the spectral controller held one.
The waveform controller used both \(\mathcal J_{\mathrm{wave}}\) and \(\mathcal J_{\mathrm{spec}}\).
These differences make the comparison one of the tuned controllers as a whole.
In the earlier reserved comparison of the spectral controllers on the same seeds, the \ac{dmdc}, \ac{mlp}, and \ac{cnn} gave similar mean burdens of \(5.04\%\), \(5.08\%\), and \(5.24\%\), at \SI{25.71}{\milli\ampere\second}, \SI{24.66} {\milli\ampere\second}, and \SI{21.92}{\milli\ampere\second}, respectively.
The available pure-\ac{mse} waveform neural runs used a tracking-only cost on five development seeds.
Their mean delivered charges fell below the \SIrange{20}{30}{\milli\ampere\second} development target. % this is irrelevant
A paired reserved comparison of charge-calibrated waveform controllers under the same cost remains open.

% Data: waveform_dmd/benchmark, observable_benchmark, and artifacts/trajopt_qp_deployment/report.md.
We assessed the condensed spectral formulation separately by replaying \num{1780} replans from the archived spectral \ac{dmdc} controller.  % weird way of testing that. Setup and run a proper experiment. Just test on the same closed-loop trajectories.
The warm online calculation, including the rollout, bounds update, and OSQP solve, took \SI{25.4}{\milli\second} at the median and \SI{34.3}{\milli\second} at the 95th percentile.
All replans in that replay completed within \SI{60}{\milli\second}, while a second replay had one overrun at \SI{68.2}{\milli\second}.
The archived IPOPT solver timings on these states had a median of \SI{1.14}{\second} and a 95th percentile of \SI{1.51}{\second}, with three unsuccessful solves. % what do you mean archived?
The replay excludes preparation, initialization, estimation, and logging, and uses states generated by the incumbent controller.
It shows that the spectral lift and the associated cost permit a much faster solve on the observed states, while closed-loop performance and end-to-end latency with the quadratic-programming backend remain to be measured.

\chapter{Waveform Controller Calibration}\label{app:waveform_controller_calibration}

\section{Training and Control Weights}\label{sec:waveform_cost_calibration}

The waveform \ac{mlp} refits in Section~\ref{sec:waveform_prediction_closed_loop} shared the architecture, training data, and initialization seed.
In the joint loss \(\mathcal{L}_\mathrm{joint}\) in~\eqref{eq:combined_predictor_loss}, the weights were \(\lambda_\mathrm{wave}=1\) and \(\lambda_\mathrm{spec}=15\) to account for the different numerical scales of the two terms.
The waveform loss \(\mathcal{L}_\mathrm{wave}\) covered the first 12 samples and the spectral loss \(\mathcal{L}_\mathrm{spec}\) covered 17 frames, following the predictability analysis in Section~\ref{sec:observable_predictability}.
The single-term fits set the other term's weight to zero.

% Protocol: docs/hyperparameter_sweep_strategy.md (Sections 6.1 and 6.3).
% Initial calibration: artifacts/sweep_2026_11/waveform_matched_geometry/mlp/effort/summary.csv.
To select a shared current weight for the suppression cost comparison, we calibrated the joint fit with both costs active.
We used \(w_{\mathrm{wave}}=1\) as the reference weight for \(\mathcal J_{\mathrm{wave}}\) and \(w_{\mathrm{spec}}=1000\) as the starting weight for \(\mathcal J_{\mathrm{spec}}\).
The waveform cost sums over channels, whereas the spectral cost averages squared excess over channels and frequency bins.
The terminal sample and frame inherited their respective stage weights.
The current weight \(w_{L_2}\) was swept over \(\{1, 3, 10, 30, 100, 300, 1000\}\) with both suppression terms active.
The selected weight, \(w_{L_2}=100\), delivered \SI{28.08}{\milli\ampere\second} on average.

% Follow-up: artifacts/waveform_joint_cost_recalibration/{report.md,rows.csv}.
% Local sensitivity and short replays: thesis/tmp/waveform_cost_weight_probe/report.md.
Holding \(w_{L_2}\) fixed while removing a suppression cost changes the balance between suppression and current effort.
To compare outcomes at similar delivered charges, we varied \(w_{\mathrm{spec}}\) and \(w_{L_2}\) with \(w_{\mathrm{wave}}=1\) (Table~\ref{tab:waveform_cost_calibration}).
The range of \(w_{\mathrm{spec}}\) reflects the relative magnitudes of the suppression costs and their sensitivity to planned currents.
For example, before \SI{2}{\second} on seed 7001 at the original weights, the median ratio of gradient norms with respect to balanced currents was 2.43 for \(w_{\mathrm{wave}}\mathcal J_{\mathrm{wave}}\) relative to \(w_{\mathrm{spec}}\mathcal J_{\mathrm{spec}}\).
A weight near 2500 would therefore bring the two local sensitivities closer on those plans.
For \(\mathcal J_{\mathrm{spec}}\) alone, we used \(w_{\mathrm{spec}}=1000\) and \(w_{L_2}=30\) to match the weight ratio of the development spectral \ac{cnn} and \ac{dmdc} settings at 17 frames.

\begin{table}
  \centering
  \caption{Suppression cost calibration for the joint waveform fit on the five development seeds.
  Burden and charge are means, and the maximum final count is the largest number of regions seizing in the final window of any run.
  The first row is the waveform tracking reference from Table~\ref{tab:waveform_loss_cost}.
  All solves were accepted.}\label{tab:waveform_cost_calibration}
  \begin{tabular}{@{}rrrrrr@{}}
    \toprule
    \(w_{\mathrm{wave}}\) & \(w_{\mathrm{spec}}\) & \(w_{L_2}\) & Burden (\%) & Charge (mAs) & Maximum final count \\
    \midrule
    1 & 0 & 100 & 5.06 & 22.93 & 4 \\
    0 & 1000 & 30 & 10.69 & 27.57 & 30 \\
    \addlinespace
    1 & 900 & 60 & 5.64 & 28.91 & 35 \\
    1 & 900 & 100 & 10.95 & 27.99 & 28 \\
    1 & 900 & 160 & 5.73 & 22.00 & 35 \\
    1 & 1500 & 60 & 5.71 & 29.81 & 35 \\
    1 & 1500 & 100 & 4.95 & 25.94 & 5 \\
    1 & 1500 & 160 & 8.98 & 24.68 & 29 \\
    1 & 2500 & 60 & 4.64 & 30.99 & 4 \\
    1 & 2500 & 100 & 4.82 & 27.32 & 4 \\
    1 & 2500 & 160 & 10.80 & 27.15 & 30 \\
    \bottomrule
  \end{tabular}
\end{table}

The timing of stimulation can matter even when total charge differs little.
In a seven-second replay of the joint fit with both suppression terms on seed 7001, reducing \(w_{L_2}\) from 100 to 10 increased charge by \SI{2.16}{\milli\ampere\second} over the replay and by \SI{2.22}{\milli\ampere\second} before \SI{2}{\second}.
The final count fell from 30 to three seizing regions.
The same weight change with only \(\mathcal J_{\mathrm{spec}}\) active on seed 7004 reduced the final count from 31 to three seizing regions.
For the waveform-only fit, the two contained runs used just \SI{5.51}{\milli\ampere\second} and \SI{6.10}{\milli\ampere\second}, while the three runs with widespread recruitment recorded 48 to 61 failed solves each.

\section{Control Horizon and Terminal Weight}

% Data: artifacts/sweep_2026_11/waveform_matched_geometry/mlp/{horizons,terminal,terminal_confirm,final_effort}.
With the joint fit and \(\mathcal J_{\mathrm{wave}}\) selected, we tested whether changing the control horizon or increasing the terminal weight improved containment at \(w_{L_2}=100\).
Increasing the terminal multiplier replaced the tracking weight on the final forecast sample while retaining weight 1 on the preceding samples.
A multiplier of 1 gave the same outcomes as the default, which inherits the stage weight.
The tested increases offered no improvement in mean burden, and the fivefold weight allowed widespread recruitment on seed 7004, which ended with 32 regions seizing.

\begin{table}
  \centering
  \caption{Development outcomes for the joint waveform \ac{mlp} with suppression cost \(\mathcal J_{\mathrm{wave}}\), averaged across five seeds.
  The terminal sweep held the horizon at 51 samples, and the horizon sweep held the terminal multiplier at 1.
  All solves were accepted.}\label{tab:waveform_horizon_terminal}
  \begin{tabular}{@{}llrrr@{}}
    \toprule
    Sweep & Setting & Burden (\%) & Charge (mAs) & Final count \\
    \midrule
    Terminal multiplier & 1 & 5.06 & 22.93 & 3.6 \\
                        & 2 & 5.13 & 22.96 & 3.6 \\
                        & 5 & 9.01 & 25.28 & 9.2 \\
                        & 10 & 5.15 & 23.57 & 4.0 \\
    \addlinespace
    Horizon & 30 samples (\SI{0.60}{\second}) & 5.17 & 19.18 & 3.6 \\
            & 42 samples (\SI{0.84}{\second}) & 5.08 & 21.82 & 3.8 \\
            & 51 samples (\SI{1.02}{\second}) & 5.06 & 22.93 & 3.6 \\
            & 60 samples (\SI{1.20}{\second}) & 5.15 & 23.41 & 4.0 \\
            & 75 samples (\SI{1.50}{\second}) & 5.08 & 23.67 & 3.8 \\
    \bottomrule
  \end{tabular}
\end{table}

Mean burden varied by only 0.11 percentage points across the horizon sweep, and all runs ended with at most five seizing regions (Table~\ref{tab:waveform_horizon_terminal}).
The shortest horizon delivered the least charge.
We retained the 51-sample horizon used in the original comparison, which spanned the same \SI{1.02}{\second} as the 17-frame spectral design and showed no improvement in mean burden from extending it.
A final current-weight sweep with this horizon and waveform tracking retained \(w_{L_2}=100\), the tested setting within the development charge range.
